% Figure for Classical Mechanics §B10.1 (energy ellipsoid omega^T I omega = 2T for principal moments I1:I2:I3 = 1:2:4, semi-axes sqrt(2T/I_k) = 2, 1.41, 1 for 2T = 4; right: its section in the (omega1, omega3) plane, omega at parameter 40 degrees, L = I omega = (1.53, 2.57) drawn at 0.5 scale, normal to the ellipse).
\begin{tikzpicture}[font=\small, >={Stealth[length=2.2mm]}]
  \begin{axis}[name=left, width=7.2cm, height=7.2cm, axis lines=none, view={35}{22}, axis equal image,
      xmin=-2.6, xmax=2.6, ymin=-2.6, ymax=2.6, zmin=-1.6, zmax=1.6, clip=false]
    \pgfplotsinvokeforeach{15,35,55,75,90,105,125,145,165}{
      \addplot3[black!30, line width=0.35pt, domain=0:360, samples=72, samples y=0]
        ({2*sin(#1)*cos(x)}, {1.4142*sin(#1)*sin(x)}, {cos(#1)});}
    \pgfplotsinvokeforeach{0,30,60,90,120,150}{
      \addplot3[black!30, line width=0.35pt, domain=0:360, samples=72, samples y=0]
        ({2*sin(x)*cos(#1)}, {1.4142*sin(x)*sin(#1)}, {cos(x)});}
    \draw[->, figB, thick] (axis cs:0,0,0) -- (axis cs:2.9,0,0) node[below left, font=\scriptsize] {$\omega_1$};
    \draw[->, figB, thick] (axis cs:0,0,0) -- (axis cs:0,2.3,0) node[right, font=\scriptsize] {$\omega_2$};
    \draw[->, figB, thick] (axis cs:0,0,0) -- (axis cs:0,0,1.6) node[above, font=\scriptsize] {$\omega_3$};
    \node[font=\scriptsize, black!70] at (axis cs:0,0,-2.3) {$I_1\omega_1^2 + I_2\omega_2^2 + I_3\omega_3^2 = 2T$};
  \end{axis}
  \begin{axis}[at={(left.east)}, anchor=west, xshift=0.6cm, width=6.4cm, axis equal image, axis lines=none,
      xmin=-2.3, xmax=2.6, ymin=-1.4, ymax=2.4, clip=false]
    \addplot[black!55, line width=0.8pt, domain=0:360, samples=120] ({2*cos(x)}, {sin(x)});
    \draw[->, black!55, thin] (axis cs:-2.3,0) -- (axis cs:2.5,0) node[right, font=\scriptsize] {$\omega_1$};
    \draw[->, black!55, thin] (axis cs:0,-1.3) -- (axis cs:0,1.6) node[above, font=\scriptsize] {$\omega_3$};
    % tangent at omega = (1.532, 0.643): direction (-2 sin t, cos t) = (-1.286, 0.766)
    \draw[black!45, dashed] (axis cs:{1.532-0.9*1.286},{0.643+0.9*0.766}) -- (axis cs:{1.532+0.5*1.286},{0.643-0.5*0.766});
    \draw[->, figB, very thick] (axis cs:0,0) -- (axis cs:1.532,0.643) node[pos=0.6, below right, font=\scriptsize] {$\vec\omega$};
    \draw[->, figR, very thick] (axis cs:0,0) -- (axis cs:0.766,1.286) ;
    \draw[->, figR, thick, dashed] (axis cs:1.532,0.643) -- (axis cs:2.298,1.929);
    \fill (axis cs:1.532,0.643) circle[radius=1.6pt];
    \node[font=\scriptsize, figR, anchor=south west] at (axis cs:0.82,1.3) {$\vec L = \mathsf I\vec\omega$};
    \node[font=\scriptsize, figR, anchor=west] at (axis cs:2.0,2.05) {normal at $\vec\omega$};
  \end{axis}
\end{tikzpicture}
